\section{En su propia historia: gestión con RL, miedo y amplificador de la civilización}

Esta sección vuelve a la perspectiva del fundador. Tim dice que hay que gestionar al estilo de RL y no solo con SFT. La metáfora resulta interesante: SFT equivale a mostrarle al equipo las respuestas correctas, mientras que RL equivale a fijar objetivos y retroalimentación para que el equipo aprenda estrategias mediante la exploración. Sin embargo, el riesgo de gestionar con RL es que resulta fácil ser objeto de hack; es decir, el equipo puede optimizar los indicadores en lugar del objetivo real.

\begin{figure}[H]
\centering
\includegraphics[width=0.96\textwidth]{rl-management-sft-risk.png}
\caption{Gestionar con RL y no solo con SFT: los objetivos, la retroalimentación y la recompensa moldean al equipo, pero también son fáciles de someter a hack. Diagrama conceptual propio, elaborado a partir de la conversación de 01:25:05--01:34:00.}
\end{figure}

\begin{knowledgebox}{Lectura de la figura: por qué es útil la metáfora de la gestión}
La gestión tipo SFT pone el acento en la demostración y en las respuestas correctas; la gestión tipo RL lo pone en los objetivos, la retroalimentación y la exploración. Tanto el entrenamiento de modelos como la gestión de organizaciones necesitan feedback, pero si el feedback está mal diseñado, termina siendo objeto de hack.
\end{knowledgebox}

\subsection{Complejidad e historia}

La subsección anterior explicó que la gestión con RL moldea al equipo; esta vuelve a cómo entiende el fundador esa complejidad. Yang Zhilin (杨植麟) dice que gran parte de la complejidad se añade de forma artificial y que, en realidad, las cosas no son tan complejas; aun así, el fundador necesita sentir, dentro de su propia historia, qué clase de persona es y por qué hace lo que hace. Esto no es huir de la técnica, sino reconocer que una empresa de modelos es la superposición de un sistema técnico y un sistema narrativo: la ruta técnica exige criterio, la organización necesita sentido y el fundador necesita seguir actuando en medio de la incertidumbre.

\begin{importantbox}{La IA es un amplificador de la civilización humana}
En la entrevista se menciona que la respuesta de Kimi a «¿qué es la IA?» es: la IA es un amplificador de la civilización humana. Esta definición eleva el modelo de la escala de herramienta a la escala de civilización: amplifica el conocimiento, la productividad y la capacidad organizativa, pero también amplifica los miedos, los malentendidos y las controversias propias de los estados intermedios.
\end{importantbox}

\begin{knowledgebox}{Tres tipos de complejidad en la historia del fundador}
\begin{tabularx}{\textwidth}{p{0.20\textwidth}p{0.34\textwidth}X}
\toprule
Complejidad & Origen & Forma de afrontarla \\
\midrule
Complejidad técnica & El modelo, los datos, la arquitectura y el producto cambian al mismo tiempo & Volver a problemas verificables y al siguiente experimento. \\
Complejidad organizativa & El equipo, la retroalimentación, la opinión pública y los indicadores de gestión se influyen entre sí & Corregir con retroalimentación, pero evitar que el reward sea objeto de hack. \\
Complejidad narrativa & Los estados intermedios pueden ser malinterpretados o criticados desde fuera & Mantener la dirección de la acción dentro de la propia historia. \\
\bottomrule
\end{tabularx}
\end{knowledgebox}

\subsection{El miedo y el paso actual}

Ante las tormentas de opinión pública y los altibajos del emprendimiento, Yang Zhilin admite que sin duda hay miedo, pero que lo más importante es concentrarse en lo que se puede hacer en el paso actual. Esta actitud se corresponde con la idea de «la montaña infinita»: si los problemas son inevitables, el miedo también lo es; pero los problemas pueden resolverse, de modo que la acción debe volver al siguiente paso.

\begin{warningbox}{Los estados intermedios siempre serán criticados}
Todo sistema complejo es imperfecto en sus estados intermedios y puede convertirse en blanco de críticas. Esto vale en especial para las empresas de modelos: las capacidades, el producto, el código abierto, la comercialización, la organización y la opinión pública se encuentran en un equilibrio dinámico. No se puede dejar de resolver problemas solo porque un estado intermedio reciba críticas.
\end{warningbox}

\subsection{Resumen de la sección}

La última sección devuelve el juicio técnico a la historia del fundador. K2 no es solo el lanzamiento de un modelo, sino también una narrativa sobre sí mismo que trata de problemas infinitos, retroalimentación organizativa, gestión del miedo y la IA como amplificador de la civilización.
